% Generated draft: requires mathematical and source review.
\documentclass[11pt]{article}
\usepackage[margin=25mm]{geometry}
\usepackage{fontspec}
\setmainfont{Latin Modern Roman}
\usepackage{amsmath,amssymb,mathtools,unicode-math}
\setmathfont{Latin Modern Math}
\usepackage{enumitem,needspace,xurl,hyperref}
\setlength{\parindent}{0pt}
\setlength{\parskip}{5pt}
\setlength{\emergencystretch}{3em}
\newcommand{\N}{\mathbb{N}}
\newcommand{\Z}{\mathbb{Z}}
\newcommand{\Q}{\mathbb{Q}}
\newcommand{\R}{\mathbb{R}}
\newcommand{\C}{\mathbb{C}}
\newcommand{\F}{\mathbb{F}}
\newcommand{\A}{\mathbb{A}}
\newcommand{\PP}{\mathbb{P}}
\newcommand{\E}{\mathbb{E}}
\newcommand{\eps}{\varepsilon}
\newcommand{\dd}{\,\mathrm{d}}
\begin{document}
\section*{UnsolvedMath ID: 30004748}
\textbf{Source problem code:} catalogue entry
\subsection{Definitions and mathematical statement}
Use the usual convention that a hyperbolic curve over \(\Q\) is a smooth, geometrically connected curve \(X/\Q\) whose smooth projective completion \(\overline X\) has genus \(g\), with \(r=\#(\overline X\setminus X)(\overline{\Q})\) satisfying \(2g-2+r>0\). The supplied record does not explicitly specify this convention.

Fix the model over \(\Z\) referred to in the source, and denote it by \(\mathcal X\), with generic fiber identified with \(X\). For a number field \(K\), let \(\mathcal O_K\) be its ring of integers and set
\[
\mathcal X(\mathcal O_K)=\operatorname{Hom}_{\Z}(\operatorname{Spec}\mathcal O_K,\mathcal X),\qquad
\mathcal X(\Z)=\operatorname{Hom}_{\Z}(\operatorname{Spec}\Z,\mathcal X).
\]
The source's notation \(X(\mathcal O_K)\) and \(X(\Z)\) is interpreted using this fixed model. These sets can depend on the model; the record does not specify additional conditions on it or clarify the quantifier over models.

Under the usual field-theoretic convention, call a number field \(K\) primitive if there is no field \(F\) with \(\Q\subsetneq F\subsetneq K\). For an integer \(n\geq1\) and a real number \(B\geq1\), define
\[
\mathcal P_n(B)=\left\{[K]:[K:\Q]=n,\ K\text{ is primitive},\ |\operatorname{Disc}(K/\Q)|\leq B\right\},
\]
where \([K]\) denotes a \(\Q\)-isomorphism class and the discriminant is the field discriminant. Thus ordering by discriminant is interpreted as counting isomorphism classes by absolute field discriminant.

\textbf{Conjecture (supplied formulation).} For every hyperbolic curve \(X/\Q\) equipped with the intended fixed model \(\mathcal X/\Z\), and every integer \(n\geq1\),
\[
\lim_{B\to\infty}
\frac{\#\left\{[K]\in\mathcal P_n(B):
\mathcal X(\mathcal O_K)=
\operatorname{im}\bigl(\mathcal X(\Z)\longrightarrow\mathcal X(\mathcal O_K)\bigr)
\right\}}
{\#\mathcal P_n(B)}=1.
\]
The map in this formula is induced by \(\Z\hookrightarrow\mathcal O_K\); equality means that every integral point over \(\mathcal O_K\) comes from an integral point over \(\Z\). Ratios are taken when the denominator is nonzero. The case \(n=1\) is included: its only field is \(\Q\), whose ring of integers is \(\Z\), so the ratio equals \(1\) for \(B\geq1\). The model qualification remains unresolved and must be checked against the source before treating this as a definitive formulation.
\subsection{Short English statement}
For a hyperbolic curve over \(\Q\) with its intended fixed integral model, the conjecture predicts that, for each degree \(n\geq1\), a proportion tending to \(100\%\) of primitive number fields, counted by absolute discriminant, introduce no integral points beyond those over \(\Z\). The precise scope of the allowed integral models requires clarification.
\subsection{Sources}
\begin{itemize}
\item Local UnsolvedMath registry, permanent ID \(30004748\), title: Integral Points over Almost All Primitive Number Fields; supplied source code \texttt{catalogue entry}, in the Oberwolfach Reports Open Problems collection. No individual proposer is identified; the supplied proposed year is \(2021\).
\item Supplied source citation: \emph{Explicit Methods in Number Theory (2022)}, Oberwolfach Reports, pp.\ 1805--1858, DOI \texttt{10.4171/owr/2021/34}. \url{https://doi.org/10.4171/owr/2021/34}. The supplied background describes these reports as documenting mathematical workshops and their open-problem discussions.
\item Supplied literature source, labeled ``Oberwolfach report'': \url{https://ems.press/content/serial-article-files/46913}.
\item Supplied literature source, labeled ``Khawaja--Siksek'': \url{https://arxiv.org/abs/2511.15635}. No title or precise result is supplied.
\item Supplied dated literature curation, checked \texttt{2026-08-22}: status labeled open; the assessment reports that no proof of the stated integral-point equality for one hundred percent of primitive number fields was found, and that recent work proves only special cases of related new-point conjectures. This is supplied dated curation derived from the OWR record, not an independently verified literature assessment. The supplied background expressly requests independent verification of time-sensitive status and cited claims.
\end{itemize}
\end{document}
